% !TEX root = ../main.tex
% ALTERNATIVA CON IMMAGINE — confrontare con 08_core_dag_cicd.tex (versione a bullet)
% NOTA: questa immagine illustra solo il DAG scheduler, uno dei tre meccanismi.
% config-driven ed exit code restano solo nello speech, non rappresentati visivamente.

\begin{frame}{Riproducibilità e integrazione CI/CD}
  \begin{adjustbox}{width=\textwidth, height=0.78\textheight, keepaspectratio, center}
    \begin{tikzpicture}[
        node distance=1.5ex and 0.8em,
        tst/.style={draw=orange!55, fill=yellow!18, rounded corners=1.5pt, minimum width=4.5em, minimum height=3.5ex, font=\footnotesize\ttfamily, align=center},
        tst-hl/.style={tst, draw=orange!70, fill=yellow!35, thick},
        ext/.style={tst, minimum width=6.5em, font=\scriptsize\ttfamily},
        phB/.style={draw=blue!40, fill=blue!7, rounded corners=1.5pt, minimum width=5.5em, minimum height=3.5ex, font=\footnotesize\ttfamily, align=center},
        phC/.style={draw=gray!40, fill=gray!8, dashed, rounded corners=1.5pt, minimum width=7em, minimum height=4.5ex, font=\footnotesize, align=center},
        bg-A/.style={draw=orange!50, fill=yellow!6, dashed, rounded corners=3pt, inner sep=12pt},
        bg-B/.style={draw=blue!40, fill=blue!4, dashed, rounded corners=3pt, inner sep=10pt},
        arr/.style={-latex, semithick, blue!60},
        darr/.style={-latex, semithick, dashed, gray!60}
      ]

      \node[tst] (d0a) {0.1};
      \node[tst, right=of d0a] (d0b) {0.2};
      \node[tst, right=of d0b] (d0c) {0.3};

      \node[tst-hl, below=of d0a] (a11) {\textbf{1.1}};
      \node[tst, right=of a11] (d1b) {1.5};
      \node[tst, right=of d1b] (d1c) {1.6};

      \node[tst, below=3.5ex of a11] (d3a) {3.3};

      \node[tst, below=of d3a] (d4a) {4.1};
      \node[tst, right=of d4a] (d4b) {4.2};
      \node[tst, right=of d4b] (d4c) {4.3};

      \node[tst, below=of d4a] (d6a) {6.2};
      \node[tst, right=of d6a] (d6b) {6.4};

      \node[tst, below=of d6a] (d7a) {7.2};

      \node[ext, below=of d7a.south west, anchor=north west] (ex0) {ext.0.1.nuclei};
      \node[ext, below=of ex0.south west, anchor=north west] (ex1) {ext.1.5.sslyze};
      \node[ext, right=of ex1] (ex2) {ext.1.5.testssl};

      \begin{scope}[on background layer]
        \node[bg-A, fit=(d0a)(d0c)(a11)(d1c)(d4c)(d7a)(ex1)(ex2)] (phaseA) {};
      \end{scope}

      \node[font=\footnotesize\bfseries, text=orange!80!black, anchor=south west]
      at ([xshift=1ex, yshift=0.5ex]phaseA.north west) {Phase A (16 test, nessuna dipendenza)};

      \path (a11.south) ++(0, -1.75ex) coordinate (bus-y);
      \path (phaseA.east |- bus-y) coordinate (exitA);
      \path (exitA) ++(2.5em, 0) coordinate (split);

      \node[phB, anchor=west] (b14) at ([xshift=1.5em, yshift=3ex]split) {1.4};
      \node[phB, anchor=west] (b21) at ([xshift=1.5em, yshift=-3ex]split) {2.1};

      \begin{scope}[on background layer]
        \node[bg-B, fit=(b14)(b21)] (phaseB) {};
      \end{scope}

      \node[font=\footnotesize\bfseries, text=blue!80!black, anchor=south west]
      at ([xshift=1ex, yshift=0.5ex]phaseB.north west) {Phase B (\texttt{depends\_on=["1.1"]})};

      \node[phC, anchor=west] (phaseC) at ([xshift=3.5em]phaseB.east |- split) {\textit{Phase C}\\{\scriptsize (pianificata)}};

      \draw[blue!60, semithick] (a11.south) |- (exitA) --
      node[above, yshift=0.3ex, font=\scriptsize\itshape, text=gray!80!black, fill=white, inner sep=0.1pt] {\texttt{1.1} $\rightarrow$} (split);

      \draw[arr] (split) |- (b14.west);
      \draw[arr] (split) |- (b21.west);

      \draw[darr] (phaseB.east |- split) --
      node[above, yshift=0.3ex, font=\scriptsize\itshape, text=gray!80!black, fill=white, inner sep=0.1pt] {dipendenze} (phaseC.west |- split);
    \end{tikzpicture}
  \end{adjustbox}
\end{frame}